\Article{21}
In den vorangehenden Rechnungen haben wir den Exponenten $n$ unbestimmt gelassen. Vom $24$.—$28$. Juni $1832$ führten wir zwei Versuchsreihen durch und dehnten sie auf die größten Abstände aus, die der Raum zuließ. Daraus wird sehr deutlich hervorgehen, welchen Wert die Natur verlangt. In der ersten Reihe wurde die zweite Nadel auf einer Geraden senkrecht zum magnetischen Meridian aufgestellt (nach der ersten Methode in Artikel $19$), in der zweiten lag der Mittelpunkt der Nadel im Meridian selbst. Die Tabelle gibt einen Überblick über die Ergebnisse dieser Versuche. Die Abstände $R$ sind in Bruchteilen eines Meters angegeben, und die Werte des Winkels $\frac14(u-u'+u''-u''')$ werden für die erste Reihe mit $v$, für die zweite mit $v'$ bezeichnet.
\begin{center}
\begin{tabular}{c|rrr|rrr}
$R$&\multicolumn{3}{c|}{$v$}&\multicolumn{3}{c}{$v'$}\\\hline
$1^{\mathrm m}1$&&&&$1^0$&$57'$&$24''8$\\
$1{,}2$&&&&$1$&$29$&$40{,}5$\\
$1{,}3$&$2^0$&$13'$&$51''2$&$1$&$10$&$19{,}3$\\
$1{,}4$&$1$&$47$&$28{,}6$&$0$&$55$&$58{,}9$\\
$1{,}5$&$1$&$27$&$19{,}1$&$0$&$45$&$14{,}3$\\
$1{,}6$&$1$&$12$&$7{,}6$&$0$&$37$&$12{,}2$\\
$1{,}7$&$1$&$0$&$9{,}9$&$0$&$30$&$57{,}9$\\
$1{,}8$&$0$&$50$&$52{,}5$&$0$&$25$&$59{,}5$\\
$1{,}9$&$0$&$43$&$21{,}8$&$0$&$22$&$9{,}2$\\
$2{,}0$&$0$&$37$&$16{,}2$&$0$&$19$&$1{,}6$\\
$2{,}1$&$0$&$32$&$4{,}6$&$0$&$16$&$24{,}7$\\
$2{,}5$&$0$&$18$&$51{,}9$&$0$&$9$&$36{,}1$\\
$3{,}0$&$0$&$11$&$0{,}7$&$0$&$5$&$33{,}7$\\
$3{,}5$&$0$&$6$&$56{,}9$&$0$&$3$&$28{,}9$\\
$4{,}0$&$0$&$4$&$35{,}9$&$0$&$2$&$22{,}2$
\end{tabular}
\end{center}

Schon ein flüchtiger Blick auf diese Zahlen zeigt, dass bei größeren Werten die Zahlen der zweiten Spalte ungefähr doppelt so groß sind wie die der dritten und die Zahlen beider Spalten ungefähr im umgekehrten Verhältnis zum Kubus der Abstände stehen. Somit kann kein Zweifel daran bleiben, dass der Wert $n=2$ richtig ist. Um dieses \SourceLeaf{114}Gesetz in den einzelnen Versuchen noch weiter zu bestätigen, haben wir alle Zahlen nach der Methode der kleinsten Quadrate behandelt und die folgenden Koeffizienten erhalten.
\[
\begin{aligned}
\operatorname{tang}v&=0{,}086870R^{-3}-0{,}002185R^{-5}\\
\operatorname{tang}v'&=0{,}043435R^{-3}+0{,}002449R^{-5}
\end{aligned}
\]

Hier folgt der Vergleich der nach diesen Formeln berechneten Werte mit den beobachteten Werten.
\begin{center}
\textit{Berechnete Werte.}\par
\nopagebreak[4]
\setlength{\tabcolsep}{4pt}
\begin{tabular}{c||rrr|r||rrr|r}
$R$&\multicolumn{3}{c|}{$v$}&Differenz&\multicolumn{3}{c|}{$v'$}&Differenz\\\hline
$1^{\mathrm m}1$&&&&&$1^0$&$57'$&$22''0$&$+2''8$\\
$1{,}2$&&&&&$1$&$29$&$46{,}5$&$-6{,}0$\\
$1{,}3$&$2^0$&$13'$&$50''4$&$+0''8$&$1$&$10$&$13{,}3$&$+6{,}0$\\
$1{,}4$&$1$&$47$&$24{,}1$&$+4{,}5$&$0$&$55$&$58{,}7$&$+0{,}2$\\
$1{,}5$&$1$&$27$&$28{,}7$&$-9{,}6$&$0$&$45$&$20{,}9$&$-6{,}6$\\
$1{,}6$&$1$&$12$&$10{,}9$&$-3{,}3$&$0$&$37$&$15{,}4$&$-3{,}2$\\
$1{,}7$&$1$&$0$&$14{,}9$&$-5.0$&$0$&$30$&$59{,}1$&$-1{,}2$\\
$1{,}8$&$0$&$50$&$48{,}3$&$+4{,}2$&$0$&$26$&$2{,}9$&$-3{,}4$\\
$1{,}9$&$0$&$43$&$14{,}0$&$+7{,}8$&$0$&$22$&$6{,}6$&$+2{,}6$\\
$2{,}0$&$0$&$37$&$5{,}6$&$+10{,}6$&$0$&$18$&$55{,}7$&$+5{,}9$\\
$2{,}1$&$0$&$32$&$3{,}7$&$+0{,}9$&$0$&$16$&$19{,}8$&$+4{,}9$\\
$2{,}5$&$0$&$19$&$2{,}1$&$-10{,}2$&$0$&$9$&$38{,}6$&$-2{,}5$\\
$3{,}0$&$0$&$11$&$1{,}8$&$-1{,}1$&$0$&$5$&$33{,}9$&$-0{,}2$\\
$3{,}5$&$0$&$6$&$57{,}1$&$-0{,}2$&$0$&$3$&$29{,}8$&$-1{,}0$\\
$4{,}0$&$0$&$4$&$39.6$&$-3{,}7$&$0$&$2$&$20{,}5$&$+1{,}7$
\end{tabular}
\end{center}

